\chapter{ER Diagram}

แผนภาพความสัมพันธ์ระหว่างข้อมูล (Entity Relationship Diagram) ของระบบบริหารจัดการยิม (GymBro Management System) แสดงถึงโครงสร้างการจัดเก็บข้อมูลเชิงสัมพันธ์ (Relational Database) โดยเน้นที่การเชื่อมโยงข้อมูลหลัก 4 ส่วน ได้แก่ ข้อมูลลูกค้า (Customers), ข้อมูลเทรนเนอร์ (Trainers), ข้อมูลอุปกรณ์ (GymEquipments), และข้อมูลการจอง (TrainingSessions) ซึ่งมีความสัมพันธ์กันดังปรากฏในรูปภาพต่อไปนี้

\begin{figure}[h]
\centering
\includegraphics[width=0.9\textwidth]{Image/er_diagram.png}
\caption{ER Diagram แสดงความสัมพันธ์ระหว่างตารางในฐานข้อมูล}
\label{fig:er_diagram}
\end{figure}

\section{คำอธิบายความสัมพันธ์ (Relationships Explanation)}
ระบบฐานข้อมูลถูกออกแบบตามหลักการ Normalization เพื่อลดความซ้ำซ้อนของข้อมูล โดยมีความสัมพันธ์แบบ One-to-Many (1:N) ระหว่างตารางหลัก (Master Tables) และตารางรายการเคลื่อนไหว (Transaction Table) ดังนี้:

\begin{enumerate}
    \item \textbf{Customers (1) $\rightarrow$ TrainingSessions (N)}: 
    ความสัมพันธ์ระหว่างลูกค้าและการจอง เป็นแบบ 1-ต่อ-กลุ่ม โดยลูกค้าหนึ่งคน (Customer) สามารถทำรายการจองเซสชันการฝึกซ้อม (Training Sessions) ได้หลายครั้งในช่วงเวลาที่ต่างกัน แต่แต่ละรายการจองจะเป็นของลูกค้าเพียงคนเดียว
    
    \item \textbf{Trainers (1) $\rightarrow$ TrainingSessions (N)}: 
    ความสัมพันธ์ระหว่างเทรนเนอร์และการจอง เป็นแบบ 1-ต่อ-กลุ่ม โดยเทรนเนอร์หนึ่งคน (Trainer) สามารถรับผิดชอบดูแลการฝึกซ้อมได้หลายเซสชัน แต่ในหนึ่งเซสชันจะมีเทรนเนอร์ผู้ดูแลหลักเพียงคนเดียว
    
    \item \textbf{GymEquipments (1) $\rightarrow$ TrainingSessions (N)}: 
    ความสัมพันธ์ระหว่างอุปกรณ์และการจอง เป็นแบบ 1-ต่อ-กลุ่ม โดยอุปกรณ์หนึ่งชิ้น (Equipment) สามารถถูกจองใช้งานได้ในหลายเซสชัน (คนละช่วงเวลา) แต่ในแต่ละเซสชันจะมีการระบุการจองใช้อุปกรณ์หลักเพียงหนึ่งชิ้น
\end{enumerate}

\section{รายละเอียดโครงสร้างตาราง (Table Schema Detail)}
รายละเอียดของคอลัมน์ (Attributes) ประเภทข้อมูล (Data Types) และข้อกำหนด (Constraints) ของแต่ละตาราง มีดังนี้:

\subsection{1. ตาราง Customers (ข้อมูลสมาชิกและโปรไฟล์ผู้ใช้)}
ตารางหลักสำหรับเก็บข้อมูลส่วนตัวของผู้ใช้งาน (User Profile) และสถานะสมาชิก
\begin{itemize}
    \item \textbf{id} (PK): รหัสสมาชิก (Integer, Auto Increment, SERIAL) เป็น Primary Key สำหรับอ้างอิง
    \item \textbf{fullName}: ชื่อ-นามสกุลจริงของผู้ใช้งาน (VARCHAR(150), NOT NULL)
    \item \textbf{email}: อีเมลสำหรับเข้าสู่ระบบและติดต่อ (VARCHAR(150), NOT NULL, UNIQUE) ห้ามซ้ำกันในระบบ
    \item \textbf{phone}: เบอร์โทรศัพท์ (VARCHAR(20)) ใช้เป็นรหัสผ่านทางเลือกในการเข้าสู่ระบบ
    \item \textbf{memberType}: ประเภทระดับสมาชิก (VARCHAR(50), NOT NULL) เช่น 'Standard', 'Premium', 'VIP'
    \item \textbf{memberLevel}: ระดับความสามารถสมาชิก (VARCHAR(50), NOT NULL) เช่น 'Basic', 'Silver', 'Gold'
    \item \textbf{memberStartDate}: วันที่และเวลาเริ่มต้นเป็นสมาชิก (TIMESTAMP, NOT NULL)
    \item \textbf{memberEndDate}: วันที่และเวลาสิ้นสุดสมาชิก (TIMESTAMP) หากเป็น NULL หมายถึงสมาชิกตลอดชีพ
\end{itemize}

\subsection{2. ตาราง Trainers (ข้อมูลเทรนเนอร์)}
ตารางเก็บข้อมูลรายชื่อเทรนเนอร์และคุณสมบัติ
\begin{itemize}
    \item \textbf{id} (PK): รหัสเทรนเนอร์ (Integer, Auto Increment, SERIAL)
    \item \textbf{trainerName}: ชื่อเทรนเนอร์ (VARCHAR(150), NOT NULL)
    \item \textbf{trainerLevel}: ระดับประสบการณ์ของเทรนเนอร์ (VARCHAR(50), NOT NULL) เช่น 'Junior', 'Senior', 'Master'
    \item \textbf{specialty}: ความเชี่ยวชาญพิเศษ (VARCHAR(100)) เช่น 'Yoga', 'Body Building', 'Cardio'
    \item \textbf{phone}: เบอร์โทรศัพท์ติดต่อ (VARCHAR(20))
\end{itemize}

\subsection{3. ตาราง GymEquipments (ข้อมูลอุปกรณ์)}
ตารางเก็บข้อมูลอุปกรณ์ออกกำลังกายที่ให้บริการในยิม
\begin{itemize}
    \item \textbf{id} (PK): รหัสอุปกรณ์ (Integer, Auto Increment, SERIAL)
    \item \textbf{equipmentName}: ชื่อเรียกอุปกรณ์ (VARCHAR(150), NOT NULL)
    \item \textbf{category}: หมวดหมู่อุปกรณ์ (VARCHAR(100)) เช่น 'Strength', 'Cardio', 'Flexibility'
    \item \textbf{status}: สถานะความพร้อมใช้งาน (VARCHAR(50), NOT NULL) เช่น 'Available', 'In Use', 'Maintenance'
\end{itemize}

\subsection{4. ตาราง TrainingSessions (ตารางการจองฝึกซ้อม)}
ตารางเก็บข้อมูล Transaction การจองคลาสหรือเซสชันการฝึกซ้อม เชื่อมโยงความสัมพันธ์ (Foreign Keys) ไปยังตารางอื่นๆ
\begin{itemize}
    \item \textbf{id} (PK): รหัสรายการจอง (Integer, Auto Increment, SERIAL)
    \item \textbf{sessionDate}: วันและเวลาเริ่มต้นการฝึกซ้อม (TIMESTAMP, NOT NULL)
    \item \textbf{endDate}: วันและเวลาสิ้นสุดการฝึกซ้อม (TIMESTAMP) คำนวณอัตโนมัติจากระยะเวลาเซสชัน
    \item \textbf{customerId} (FK): รหัสสมาชิกผู้จอง (Integer, NOT NULL) \newline
    อ้างอิง (REFERENCES) ไปยังตาราง \textbf{Customers(id)} \newline
    กำหนด \textbf{ON DELETE CASCADE} (หากลบสมาชิก ประวัติการจองจะถูกลบด้วย)
    \item \textbf{trainerId} (FK): รหัสเทรนเนอร์ผู้ดูแล (Integer, NOT NULL) \newline
    อ้างอิง (REFERENCES) ไปยังตาราง \textbf{Trainers(id)}
    \item \textbf{equipmentId} (FK): รหัสอุปกรณ์ที่ใช้ (Integer, NOT NULL) \newline
    อ้างอิง (REFERENCES) ไปยังตาราง \textbf{GymEquipments(id)}
\end{itemize}
